\section{Introduction}\label{sec:introduction}

Let $A_0/\Q$ be an elliptic curve and let $p$ be a rational prime.
Write $\Sha(A_0/\Q)$ for the kernel of localization from the
Weil--Ch\^atelet group to the product of the local Weil--Ch\^atelet
groups. The full $p$-power Selmer group, with local Kummer conditions
at every place, fits into the exact sequence
\[
 0\longrightarrow A_0(\Q)\otimes\Q_p/\Z_p
 \longrightarrow\Sel_{p^\infty}(A_0/\Q)
 \longrightarrow\Sha(A_0/\Q)[p^\infty]
 \longrightarrow0.
\]
Thus its corank measures both the Mordell--Weil rank and any divisible
$p$-primary contribution from $\Sha$. Put
\[
 s_p(A_0)=\corank_{\Z_p}\Sel_{p^\infty}(A_0/\Q),
 \qquad a(A_0)=\ord_{s=1}L(A_0,s).
\]
The Selmer converse asks whether this arithmetic invariant at a
specified prime determines the analytic order of vanishing. We prove
the converse in coranks zero and one without restrictions on the curve
or the prime.

\begin{theorem}[The unrestricted low-corank Selmer converse]
\label{thm:main}
For every elliptic curve $A_0/\Q$, every prime $p$, and
$r\in\{0,1\}$,
\[
 s_p(A_0)=r
 \quad\Longrightarrow\quad
 a(A_0)=\rank_{\Z}A_0(\Q)=r,
 \qquad \#\Sha(A_0/\Q)<\infty.
\]
There is no restriction on reduction type, complex multiplication,
the residual representation, rational torsion, or rational isogenies.
\end{theorem}

The finiteness conclusion concerns the entire Tate--Shafarevich
group, not just its $p$-primary part. In particular, the Kummer sequence
then shows that all prime-power Selmer coranks agree with the rank
determined at $p$. The theorem gives rank equality and finiteness; it
does not evaluate the leading coefficient in the refined
Birch--Swinnerton-Dyer formula.

A concrete application comes from Sylvester's cube-sum problem:
which rational primes are sums of two rational cubes?
Section~\ref{sec:sylvester} applies the theorem at $p=3$ to show that,
for every prime $\ell\equiv4,7,8\pmod9$, the cubic
$X^3+Y^3=\ell Z^3$ has analytic and Mordell--Weil rank one and finite
$\Sha$. These cases have also been treated by direct Heegner-point
methods, including recent work of Yin and Burungale--Tian
\citep{Yin2026Sylvester47,Yin2026SylvesterGZ,BurungaleTian2026Sylvester}.
The deduction here illustrates the use of an unrestricted converse at
a prime of additive reduction: Satg\'e's paired isogeny descents,
with their Tate--Shafarevich terms retained, bound the full
$3$-power Selmer corank by one, and the negative root sign selects
analytic rank one.

\subsection{The converse problem and its development}

The conjecture of Birch and Swinnerton-Dyer grew from computations
relating rational points to the behavior of the elliptic $L$-function
near its central point \citep{BirchSwinnertonDyer1965}.
Coates and Wiles proved an early analytic-to-arithmetic implication for
elliptic curves with complex multiplication \citep{CoatesWiles1977}.
Modularity supplies analytic continuation and the functional equation
for every elliptic curve over $\Q$ \citep{BCDT2001}. The height formula
of Gross and Zagier and Kolyvagin's Euler systems give rank equality
and finiteness of $\Sha$ when the analytic rank is at most one
\citep{GrossZagier1986,Kolyvagin1990}. Kato's zeta elements provide
another fundamental route from a nonzero central value to arithmetic
finiteness \citep{Kato2004}.
The converse reverses this direction: a small Selmer group must force
a nonzero central value or derivative. More generally, the
$p$-converse predicts $a(A_0)=s_p(A_0)$; see
\citet[Conjecture~A]{KellerYin2024}.

In rank zero, Skinner--Urban's main-conjecture theorem gives the
Iwasawa divisibility opposite to the one obtained from Kato's Euler
system. It yields the converse at odd good ordinary primes under
residual irreducibility and an auxiliary ramification hypothesis
\citep{SkinnerUrban2014}.
In rank one, the problem becomes the non-torsion of a Heegner point.
Skinner uses anticyclotomic Iwasawa theory to deduce this from a
one-dimensional Selmer group with nonzero localization at $p$, under
ordinary, residual, and ramification hypotheses
\citep[Theorem~B]{Skinner2020}.
Wei Zhang proves nonvanishing of derived Heegner classes by
level-raising and rank-lowering arguments. Under his ordinary and
residual hypotheses, this proves the corank-one converse without the
nonzero-localization assumption \citep[Theorem~1.4]{Zhang2014}.
These two approaches have different hypotheses; their common task is
to recover analytic nonvanishing from arithmetic information that
does not itself exhibit a non-torsion Heegner point.

Later work has removed many of the restrictions in these arguments.
For good ordinary $p>3$ with irreducible $A_0[p]$,
Burungale--Castella--Skinner remove the earlier auxiliary
residual-ramification conditions \citep{BCS2025}.
Reducible residual representations require a different analysis:
Castella--Grossi--Lee--Skinner prove the converse at odd good
Eisenstein primes with a local-character restriction, and Keller--Yin
remove that restriction and treat multiplicative and potentially good
ordinary cases, including rational $p$-torsion and some additive
reduction \citep{CGLS2022,KellerYinSemistable2024,KellerYin2024}.
At good supersingular primes greater than three, Castella--Wan prove
the corank-one converse for semistable curves \citep{CastellaWan2024}.
Castella treats multiplicative reduction under local and residual
hypotheses \citep[Theorems~1.1 and~1.3]{Castella2024}; the zeta-element
constructions of Burungale--Skinner--Tian--Wan yield further
supersingular rank-zero and ordinary applications
\citep[Theorems~1.6 and~1.10]{BSTW2024}.

Complex multiplication also permits small-prime results.
Burungale--Castella--Skinner--Tian prove the corank-one converse at
good ordinary primes, including $2$ and $3$, when the associated
Hecke character has conductor exactly divisible by the different of
the CM field \citep[Theorem~A]{BCST2022}.
Burungale--Tian prove the rank-zero converse for every CM elliptic
curve at every prime \citep{BurungaleTian2026}.
Kriz's preprint states a low-corank converse at the ramified prime
for curves with CM by the full ring of integers, together with its
cube-sum application \citep[Theorems~1.1 and~1.3]{Kriz2022}.

The present proof uses two distinct results of
\citet{OpenAIGoldfeldDensity2026}: the unrestricted converse at $2$
settles that prime, while the analytic twist-density theorem supplies
the auxiliary fields for the odd-prime argument. Their precise
statements are recorded in Section~\ref{sec:reduction}.
The remaining construction keeps the specified odd prime arbitrary.
In particular, it must control local conditions even when the
elliptic curve is not ordinary and the residual representation is
reducible.

\subsection{Proof strategy: two incompatible lengths}

Fix an odd prime $p$ and suppose $s_p(A_0)\in\{0,1\}$.
The first goal is to produce a non-torsion Heegner point. The density
input selects an imaginary quadratic field $K$ and a biquadratic CM
field $L$ containing $K$. The rational Selmer spaces of $A_0$ and an
auxiliary quadratic twist over $K$ are both lines; the latter is
generated by a known non-torsion point.
Section~\ref{sec:reduction} shows that non-torsion
of the conductor-one Heegner trace $y$ implies the desired conclusion
by Gross--Zagier and Kolyvagin. Assume, for a contradiction, that $y$
is torsion. We will obtain more Galois extensions than the Selmer
hypothesis allows.

\paragraph{The cohomological upper bound.}
Let $V=T_pA_0\otimes_{\Z_p}\Q_p$. At a sequence of split auxiliary
primes $q_i$, cyclic $p$-quotients of growing ring class groups give
a tame character $\psi$ with coefficients in
\[
 R_b=k[t]/(t^b),\qquad b\ge1,
\]
where $k/\Q_p$ is finite and $\psi\equiv1\pmod t$.
We form $\psi$ by taking the $p$-adic
ultralimit of the finite-stage characters on sequences of Galois
elements over $K$, and restrict it to $L$ below.
Its cohomology is represented by
bounded finite-stage cochains whose identities hold to uniformly
increasing precision; Section~\ref{sec:tame} calls these cochains
\emph{admissible}.
Choose a place $v$ of $K$ above $p$ and the places of $L$ above $v$.
The one-sided Greenberg group
$\mathcal H_b=\GrH(L,V\otimes\psi)$ allows unrestricted local
classes at the chosen places and requires zero classes at their
conjugates, with relaxed conditions at the remaining fixed places.

If the Selmer line of $A_0/K$ has nonzero localization at $p$, then
$\mathcal H_b=0$. Otherwise that line is everywhere locally zero;
we call this the strict case. There a Heisenberg commutator modifies
the tame direction so
that no nonzero base Greenberg class lifts to first order. The exact
coefficient sequence then confines every admissible class to the
last $t$-layer and gives $\length_{R_b}\mathcal H_b\le2$, independently
of $b$. Consequently, with
\[
 b_*=1\quad\text{outside the strict case},\qquad
 b_*=3\quad\text{in the strict case},
\]
the cohomological upper bound is strictly smaller than $b_*$.

\paragraph{A theta congruence retaining the tame parameter.}
The weighted Heegner logarithm is a truncated $t$-series formed from
unnormalized sums of local logarithms of ring-class Heegner points, weighted
by the finite-stage tame characters. The torsion assumption on $y$
makes its constant coefficient tend to zero. In the strict case,
conjugation and the zero localizations give the same conclusion for
the next two coefficients. This logarithmic vanishing is distinct
from the first-order obstruction used for the upper bound. It is
the input to the automorphic construction.

Sections~\ref{sec:theta}--\ref{sec:content} construct a holomorphic
theta family on a unitary group of signature $(3,1)$ at both real
places of the maximal real subfield of $L$. The construction uses
Weil's oscillator representation and the theta and period framework
of Kudla and Rallis \citep{Weil1964,Kudla1994,Rallis1984Injectivity}.
An algebraic square identity relates its constant term to toric
periods, one of which specializes to the Heegner logarithm. The
passage from raised CM values through an inverse differential
operator to that logarithm follows the method of
Bertolini--Darmon--Prasanna \citep{BertoliniDarmonPrasanna2013}; the
argument here includes the additional bad-reduction comparisons.
Uniform integral moment estimates make the constant terms small,
while a separate positive Fourier--Jacobi coefficient has bounded
divisibility and prevents the entire theta series from disappearing.

The geometric input to cusp lifting is the ordinary locus of the
target PEL variety, not ordinary reduction of $A_0$.
Using the compactification and ordinary deformation frameworks of
Lan and Katz \citep{Lan2013,Lan2018Ordinary,Katz1981SerreTate},
Section~\ref{sec:cusp-lifting} lifts the theta series to genuine
cusp forms modulo increasing powers of $p$. Its integral coefficient
lattice and normalized Hecke traces preserve the whole truncated
tame ring. Evaluating the congruence only at individual characters
would discard the nilpotent parameter that measures the required
length.

\paragraph{The extension lower bound.}
The passage from an automorphic congruence to a Galois extension
goes back to Ribet and Wiles \citep{Ribet1976,Wiles1990}.
Skinner--Urban use congruences on larger unitary groups
\citep{SkinnerUrban2014}; the semiordinary $\mathrm{GU}(3,1)$
construction of Castella--Liu--Wan is a close antecedent allowing a
nonordinary two-dimensional source
\citep{CastellaLiuWan2022,CastellaLiuWanAddendum}.
Here Section~\ref{sec:parameters} constructs four-dimensional
characteristic-zero Galois representations from the cusp forms and
proves the local identities needed to extract one-sided Greenberg
classes. Two normalized parabolic operators are units, although
the elliptic representation remains unrestricted at $p$.

The limiting trace separates into a two-dimensional twist of $V$
and two characters; this is not a decomposition assumed for each
cuspidal Galois representation. Section~\ref{sec:extensions} uses
the off-diagonal modules of a finite limiting matrix algebra to
record extensions between these constituents, in the spirit of
the generalized matrix algebras of Bella\"iche--Chenevier
\citep{BellaicheChenevier2009}.
The main tasks are to eliminate an unwanted cyclotomic extension
quotient between the two characters and to impose the strict local
condition on extensions
of one character by the two-dimensional constituent. Uniformly
many algebra generators reduce each limiting functional to finitely
many matrix values, so the resulting cochains remain admissible.
The local identities hold over the entire Artin ring, and a Fitting
ideal calculation gives
\[
 \length_{R_{b_*}}\mathcal H_{b_*}\ge b_*.
\]
This contradicts the upper bound and proves that $y$ is non-torsion.
Figure~\ref{fig:proof} summarizes the comparison.

\begin{figure}[tb]
\centering
\begin{tikzpicture}[
  box/.style={draw=black!55,rounded corners=2pt,align=center,
    text width=5.4cm,minimum height=1.05cm,inner sep=6pt,font=\small},
  arr/.style={-{Latex[length=2mm]},thick,draw=black!65},
  node distance=0.65cm and 0.55cm]
 \node[box,text width=12cm] (hyp) at (0,0)
 {Assume the conductor-one Heegner point is torsion.\\
 Choose the tame character over $R_{b_*}$.};
 \node[box] (upper) at (-3.15,-1.9)
 {Cohomological upper bound\\
 (\cref{prop:green-length})\\
 $\length_{R_{b_*}}\GrH<b_*$};
 \node[box] (theta) at (3.15,-1.9)
 {Theta estimates and\\
 full-ring cusp congruence\\
 (\cref{prop:cusp-congruence})};
 \node[box] (lower) at (3.15,-3.8)
 {Galois extension extraction\\
 (\cref{prop:extension-extraction})\\
 $\length_{R_{b_*}}\GrH\ge b_*$};
 \node[box] (end) at (-3.15,-3.8)
 {Contradiction in both cases\\
 The Heegner point is non-torsion.};
 \draw[arr] (hyp.south -| upper.north) -- (upper.north);
 \draw[arr] (hyp.south -| theta.north) -- (theta.north);
 \draw[arr] (theta) -- (lower);
 \draw[arr] (upper) -- (end);
 \draw[arr] (lower) -- (end);
\end{tikzpicture}
\caption{The two bounds concern the same admissible Greenberg module
 over $R_{b_*}$. Nonzero localization gives upper bound $0$ and
 $b_*=1$; the strict case gives upper bound at most $2$ and $b_*=3$.}
\label{fig:proof}
\end{figure}


The proof therefore separates into the arithmetic reduction and
upper bound in Sections~\ref{sec:reduction}--\ref{sec:tame}, the theta
and cusp congruence in Sections~\ref{sec:theta}--\ref{sec:cusp-lifting},
and the Galois construction and extraction in
Sections~\ref{sec:parameters}--\ref{sec:extensions}.
Section~\ref{sec:sylvester} gives the cube-sum application.
The simultaneous transfer in Section~\ref{subsec:simultaneous-transfer}
uses Labesse's base-change theorem, with its sign property verified
for the holomorphic discrete series at both real places.
Appendix~\ref{sec:nonstandard-transfer} records a separate normalized
global measure comparison and a change-of-characteristic principle
for nonstandard weighted orbital integrals.
